\pdfbookmark[1]{Solution of Quadratic Equations. 1840.}{quadratisch1840}\par\Needspace{10\baselineskip}\begin{center}\rule{0.7\linewidth}{0.4pt}\par \textsc{Vega} and \textsc{Hülsse.} Sammlung mathematischer Tafeln. 1840.\par\rule{0.7\linewidth}{0.4pt}\end{center}

\noindent\textit{Solution of quadratic equations in the form in which the coefficients}
\textit{of the equation themselves are not given, but their logarithms, and in which one needs}
\textit{not the roots themselves (or one of them), but rather their logarithms}
\textit{for further use.}

Carrying out this task with the ordinary logarithm tables alone necessarily requires four lookups. Two lookups suffice if one uses either the trigonometric logarithm tables or the tables for logarithms of sums and \SourceLeaf{265} differences. A single lookup suffices if the latter table is extended, as here, by the additional columns $D$, $E$ and $F$.

The fourth column, $D$, contains $B+C$, the fifth, $E$, contains $A+C$,
the sixth, $F$, contains $B-A$. If we suppose that the numbers themselves to which the logarithms $A$, $B$, $C$ belong are respectively $a$, $b$, $c$, then the fourth column would contain the logarithms of $bc$, the fifth those of $ac$, the sixth those of $\frac{b}{a}$.

It may also be pointed out that the logarithms in the sixth column $F$, which are positive up to $A=0{,}208$, become negative from $A=0{,}209$
onwards (p.\ 640). It makes little difference whether they are given in this negative form or, as here, by their complements; thus, for example, for $A=0{,}367$ under $F$ one could have $-0{,}21180$ or, as here, $9{,}78820$.

To apply these additional columns to the solution of the quadratic equation
\[
pxx+qx+r=0
\]
itself, four different cases must be distinguished:
\begin{enumerate}
\item[I.] $p$ and $r$ have the same signs and $\frac{qq}{pr}$ is not less than $4$;
\item[II.] $p$ and $r$ have the same signs and $\frac{qq}{pr}$ is less than $4$;
\item[III.] $p$ and $r$ have opposite signs and $-\frac{pr}{qq}$ is greater than $2$;
\item[IV.] $p$ and $r$ have opposite signs and $-\frac{pr}{qq}$ is less than $2$.
\end{enumerate}
The case in which $-\frac{pr}{qq}=2$ holds can be assigned either to III.\ or to IV.

\Needspace{24\baselineskip}
In case II.\ the roots are imaginary; in the other cases each of the two roots is obtained in two ways. The following scheme makes everything easy to take in. Here
\[
\frac{q}{p}=h,\quad \frac{r}{q}=g
\]
has been set, partly for brevity and partly because the calculation is indeed most conveniently carried out in this form.

\begin{center}
\begingroup\renewcommand{\arraystretch}{1.7}\begin{tabular}{rlll}
 & & \emph{First root.} & \emph{Second root.} \\
I. & $+\frac{h}{g}=bc=d$ & $-\frac{h}{b}=-gc$ & $-gb=-\frac{h}{c}$ \\
II. & $+\frac{h}{g}<4$ & Imaginary. & Imaginary. \\
III. & $-\frac{g}{h}=ac=e$ & $+ha=-\frac{g}{c}$ & $+\frac{g}{a}=-hc$ \\
IV. & $-\frac{g}{h}=\frac{b}{a}=f$ & $+\frac{h}{a}=-\frac{g}{b}$ & $+ga=-hb$
\end{tabular}\endgroup
\end{center}
\par
\SourceLeaf{266}
Each reader will easily be able to work out the proofs of these rules. For good measure, an example of their use may be added.

\begin{center}
\begin{tabular}{rl}
Let the given data be & $\log p=0{,}69897$ \\
 & $\log q=0{,}84510$ \\
 & $\log r=0{,}77815\,\mathrm{n}$ \\
Thus & $\log h=0{,}14613$ \\
 & $\log g=9{,}93305\,\mathrm{n}$ \\
 & $\log(-\frac{g}{h})=9{,}78692$
\end{tabular}
\end{center}

One sees immediately that case IV.\ applies here, so that $9{,}78692$ must be sought in the sixth column under $B-A$ or $F$. The only part of the table (p.\ 643) needed here would be the following:
\begin{center}
\begin{tabular}{c|c|c}
$A$ & $B$ & $B-A$ or $F$ \\
$0{,}367$ & $0{,}15520$ & $9{,}78820$ \\
$0{,}368$ & $0{,}15489$ & $9{,}78689$
\end{tabular}
\end{center}
Thus $F=9{,}78692$ corresponds to
\begin{center}
\begingroup\setlength{\tabcolsep}{3pt}\begin{tabular}{l}
$A=0{,}36798=\log a$, from which $\log\frac{h}{a}=9{,}77815$, $\log ga=0{,}30103\,\mathrm{n}$ or \\
$B=0{,}15490=\log b$, from which $\log(-\frac{g}{b})=9{,}77815$, $\log(-hb)=0{,}30103\,\mathrm{n}$.
\end{tabular}\endgroup
\end{center}

